% Part: first-order-logic
% Chapter: axiomatic-deduction
% Section: provability-quantifiers

% verification of properties of provability needed for maximally
% consistent sets in the completeness chapter.

\documentclass[../../../include/open-logic-section]{subfiles}

\begin{document}

\olfileid{fol}{axd}{qpr}

\olsection{\usetoken{S}{derivability} او کميت ټاکونکي}

\begin{explain}
  د بشپړوالي قضيه دا هم غواړي چې بديهي استنتاجونه د~$\Proves$ په
  اړه په دې برخه کښې ثابت شوي حقيقتونه ورکړي۔
\end{explain}

\begin{thm}
\ollabel{thm:strong-generalization} که $c$ داسې !!a{constant} وي چې
په~$\Gamma$ يا~$!A(x)$ کښې نۀ راځي، او $\Gamma \Proves !A(c)$ وي،
نو $\Gamma \Proves \lforall[x][!A(x)]$۔
\end{thm}

\begin{proof}
د استنتاج د قضیې له مخې $\Gamma \Proves \ltrue \lif !A(c)$۔ ځکه
$c$ په~$\Gamma$ يا~$\top$ کښې نۀ راځي، ترلاسه کوو چې
$\Gamma \Proves \ltrue \lif \lforall[x][!A(x)]$۔ ځکه رښتيا يو
بديهي اصل دے، د ميتا-مودوس پوننس له مخې
$\Gamma \Proves \lforall[x][!A(x)]$۔
\textbf{د ژباړې د نسخې سمون (OLAX-009):} سرچينې وروستے ګام د
استنتاج قضیې ته يو ځل بيا منسوب کړے و؛ هغه قضيه يوازې د رښتيا فرض
ورزياتوي، حال دا چې دلته رښتيا بديهي اصل او مودوس پوننس فرض بېرته
لرې کوي۔
\end{proof}

\begin{prop}
\ollabel{prop:provability-quantifiers}
\begin{tagenumerate}{prvEx,prvAll}
\tagitem{prvEx}{$!A(t) \Proves \lexists[x][!A(x)]$.}{}

\tagitem{prvAll}{$\lforall[x][!A(x)] \Proves !A(t)$.}{}
\end{tagenumerate}
\end{prop}

\begin{proof}
\begin{tagenumerate}{prvEx,prvAll}
\tagitem{prvEx}{د \olref[qua]{ax:q2} او د استنتاج د قضیې له مخې۔}{}
\tagitem{prvAll}{د \olref[qua]{ax:q1} او د استنتاج د قضیې له مخې۔}{}
\end{tagenumerate}
\end{proof}

\end{document}
